\section{The Knit}

One law moves the whole mesh, and it is the same law everywhere and at every beat. It is called the \emph{knit}, because it interlaces the tones meeting in a dock and then advances them a row at a time, the way knitting works a fabric. A single beat is two steps. First the tones inside each dock \emph{collide}, a local rewrite that touches that dock alone. Then every tone \emph{streams}, moving one dock along the site it points down. The collide is where anything interesting happens, and the stream is mere transport, but together they are all there is.

A collision has to obey the constraints of the base, so it is a permutation of a dock's states that is local, reversible, and charge-conserving. The committed knit gets this with maximum economy, by pairing the sites. Each of a dock's twenty-four sites has an opposite, the direction pointing the other way, and a direction together with its opposite is a \emph{line}, the channel along which two tones pass head-on through the dock. There are twelve such lines, and the collision runs one small fixed table on each of them, independently. A line holds two tones, so it has $3 \times 3 = 9$ states, and the table is just a permutation of those nine.

\begin{table}[ht]
\centering
\begin{tabular}{@{}clllp{0.40\textwidth}@{}}
\toprule
\# & in & out & move & what it does \\
\midrule
1 & \texttt{(-1,-1)} & \texttt{(-1,-1)} & inert & like tones sit, nothing happens \\
2 & \texttt{(+1,+1)} & \texttt{(+1,+1)} & inert & like tones sit, nothing happens \\
3 & \texttt{(-1,~0)} & \texttt{(~0,-1)} & hop & a lone $-1$ crosses the dock \\
4 & \texttt{(~0,-1)} & \texttt{(-1,~0)} & hop & a lone $-1$ crosses the dock \\
5 & \texttt{(+1,~0)} & \texttt{(~0,+1)} & hop & a lone $+1$ crosses the dock \\
6 & \texttt{(~0,+1)} & \texttt{(+1,~0)} & hop & a lone $+1$ crosses the dock \\
7 & \texttt{(~0,~0)} & \texttt{(+1,-1)} & create & peace leans out a balanced pair \\
8 & \texttt{(+1,-1)} & \texttt{(-1,+1)} & flip & the pair reverses \\
9 & \texttt{(-1,+1)} & \texttt{(~0,~0)} & annihilate & the pair returns to peace \\
\bottomrule
\end{tabular}
\caption{The collision, written on one line as the pair $(\text{left}, \text{right})$. It runs unchanged on each of a dock's twelve lines.}
\label{tab:collision}
\end{table}

The nine states fall into three kinds of move, shown in Table~\ref{tab:collision}. Two are inert, like tones on a line simply coexisting. Four are hops, a lone charge sliding across the dock to the opposite site. The last three form a single cycle, and that cycle is the most consequential thing in the base. Peace creates a balanced pair, the pair flips, and the flipped pair annihilates back to peace, $(\,0,0) \to (+1,-1) \to (-1,+1) \to (0,0)$. Every row leaves the sum of its two tones unchanged, so the collision conserves charge exactly, and because it is a permutation of nine states it runs backward as cleanly as forward. The hops and inert states are their own inverses, and only the three-cycle has a genuine direction, which is undone by running it the other way.

Streaming is the easy half. After the collision, the tone now sitting in a site moves to the neighbouring dock that the site points at, carrying its value untouched. Nothing is created or destroyed, only relocated, so the streaming step is a bijection on the whole mesh and is exactly invertible. A free tone, one not caught in a collision, therefore travels in a straight line forever, one dock per beat, never spreading. That straight ballistic travel is what momentum is, and a tone's momentum is nothing more than which site it occupies.

\begin{result}[The beat]
A beat applies the collision to every dock at once, then streams every tone one dock along its site. Time is the integer count of beats, with nothing finer. Because both halves are invertible, the whole beat is an exact reversible permutation of the universe, and the closed bulk carries no arrow of time.
\end{result}

It is worth being clear about how much of physics is already implicit in so small a rule, none of it added by hand. Streaming alone gives propagation and momentum. The collisions give scattering, the interaction of one moving structure with another, and the energy cost of building a localized bound thing, which is the seed of mass. The create move, the one cycle that lifts a pair out of peace, makes the vacuum an active, fluctuating sea rather than an empty stage, and that activity later supplies both the photon and the force that binds a body to itself. All of this is consequence, read off the structure of collide-and-stream, not a list of separate ingredients. The full transition tables, including the discarded alternatives that helped fix this one, are collected in Appendix~\ref{app:base}.
